\documentclass{article}
\usepackage{xeCJK}
\usepackage{amsmath}
\linespread{1.5}

\title{第一次作业}

\date{2025.9}

\begin{document}

\maketitle

1.已知函数$f_1(x) = \frac{2x-1}{x+1}$,对于$n \in N^*$,定义$f_{n+1}(x) = f_1(f_n(x))$.若$f_{35}(x) = f_5(x)$,则$f_{38}(x) = $\newline
\newline
2.\text { 定义在 } \mathbf{R} \text { 上的函数 } f(x)(x \neq 1) \text { 满足 } f(x)+2 f\left(\frac{x+2002}{x-1}\right)=4015-x \text { ，则 } f(2004)=\newline
\newline
3.（2022年北京大学强基）已知$f(x)$是二次函数，$f(-2) = 0$，且$2x \leq f(x) \leq \frac{x^2+4}{2}$，则$f(10) = $\newline
\newline
4.函数$f(x) = \sqrt{\frac{x+1}{x^2+4x+7}}$的值域为\newline
\newline
5.函数$f(x) = |x-1| + |x-3| + |x-5| + |x-7|$的最小值为\newline
\newline
6.已知函数$f(x) = x^2 - 2|x| + 2$的定义域为$[a,b](a<b)$，值域为$[2a,2b]$，则符合条件的数组$(a,b)$为\newline
\newline
7.函数$f(x) = \frac{\sqrt{x^2+1}}{x-1}$的值域为\newline
\newline
8.（2022年山东预赛）已知函数$f(x)$满足，对任意实数$x,y$，有$f(xy) + f(x-y) \geq f(x+y)$,证明：对于任意实数$x$，均有$f(x) \geq 0$.\newline
\newline
\newline
\newline
\newline
\newline
9.设$a \leq 2$,求$y = (x-2)|x|$在区间$[a,2]$上的最大值和最小值.\newline
\newline
\newline
\newline
\newline
\newline
10.设$a = \log_{3}2, b = \ln2, c = 5^{-\frac{1}{2}}$，则（$\quad$）\newline
$A.a<b<c \qquad B.b<c<a \qquad c.c<a<b \qquad D.c<b<a$ \newline
\newline
11.已知函数$f(x) = |lgx|,$若$0<a<b$，且$f(a) = f(b)$，则$a+2b$的取值范围是（$\quad$）\newline
$A.(2\sqrt{2},+\infty) \qquad B.[2\sqrt{2},+\infty) \qquad c.(3,+\infty) \qquad D.[3,+\infty)$ \newline
\newline
12.方程$|x-3|^{\frac{x^4-8x+15}{x-2}} = 1$的解有（$\quad$）个\newline
$A.1 \qquad B.2 \qquad c.3 \qquad D.4$ \newline
\newline
13.方程组$\begin{cases}x^{x-y} = y^{x+y}\\y\sqrt{x} = 1 \end{cases}$的解有（$\quad$）个\newline
$A.1 \qquad B.2 \qquad c.3 \qquad D.4$ \newline
\newline
14.（2021年上海交大强基）若$a>1,b>1$且$\lg(a+b) = \lg a + \lg b$，则$\lg(a-1) + \lg(b-1) = $\newline
\newline
15.（2022年四川预赛）已知正实数$a,b$满足$a^b = (8a)^{9b}$,则$\log_{a}(4a^2)$的值为\newline
\newline
16.解方程$x^{\log_{a}x} = \frac{x^3}{a^2},x=$\newline
\newline
17.已知$a>0,b>0,\log_{9}a = \log_{12}b = \log_{16}(a+b)$,求$\frac{b}{a}$的值.
\end{document}
